\section*{अध्याय \ref{ch:b2:asexual-reproduction} --- पौधों में अलैंगिक जनन और क्लोनन}

\begin{solution}{exo:b2:asexual-reproduction:1}
\omterm{def:b2:asexual-reproduction:clone}{क्लोन}: किसी एक जनक की समसूत्री विभाजन से बनी आनुवंशिक रूप से समरूप
संतति। \omterm{def:b2:asexual-reproduction:clone}{जेनेट}: आनुवंशिक व्यक्ति, यानी वह सब जो एक ही
युग्मनज से उगा। \omterm{def:b2:asexual-reproduction:clone}{रैमेट}: किसी \omterm{def:b2:asexual-reproduction:clone}{जेनेट} की एक शरीरक्रियात्मक रूप से स्वतंत्र इकाई।
अंग: भूस्तारी (स्ट्रॉबेरी), प्रकंद (दूब), कंद
(आलू), शल्ककंद (प्याज़), मूल-प्ररोह (ऐस्पन), पत्ती के पौधक
(\emph{Kalanchoe})।
\end{solution}

\begin{solution}{exo:b2:asexual-reproduction:2}
कलिकाग्र और मूलवृंत के संवहनी कैंबियम एक पंक्ति में और संपर्क में होने
चाहिए; कैलस उस कटान पर पुल बनाता है और उसके आरपार जाइलम तथा फ्लोएम
विभेदित कर देता है। मूलवृंत जड़ें देता है (जल, खनिज, ओज,
मृदा-सहिष्णुता, आकार-नियंत्रण) और कलिकाग्र फलदार प्ररोह (यानी
क़िस्म)। हर भाग अपना ही जीनोम रखता है और कोशिकाएँ कभी संलयित नहीं होतीं:
दो \omterm{def:b2:meiosis-heredity:vocabulary}{जीन-प्ररूप} अगल-बगल, यानी एक काइमेरा; जबकि कोई संकर हर कोशिका में
दोनों जीनोमों का मिश्रण ढोता।
\end{solution}

\begin{solution}{exo:b2:asexual-reproduction:3}
\omterm{def:b2:asexual-reproduction:clone}{कायिक गुणन}: शरीर का कोई भाग जड़ पकड़कर नया पौधा बन जाता है।
\omterm{def:b2:asexual-reproduction:apomixis}{असंगजनन}: कोई \omterm{def:b2:angiosperm-reproduction:seed}{बीज} बिना निषेचन के बनता है, और उसका भ्रूण माता का
\omterm{def:b2:asexual-reproduction:clone}{क्लोन} होता है। \omterm{def:b2:asexual-reproduction:apomixis}{अनिषेकजनन}: कोई प्राणी किसी अनिषेचित
अंडाणु से विकसित होता है। इनमें से केवल \omterm{def:b2:asexual-reproduction:apomixis}{असंगजनन} \omterm{def:b2:angiosperm-reproduction:seed}{बीज} रखता है --- उसके
आवरण, \omterm{def:b2:angiosperm-reproduction:seed}{प्रसुप्ति} और प्रकीर्णन सहित।
\end{solution}

\begin{solution}{exo:b2:asexual-reproduction:4}
बहिर्भाग को निर्जीवाणुकृत कीजिए; उसे शर्करा, खनिज,
विटामिन तथा संतुलित ऑक्सिन और साइटोकाइनिन वाले अगार पर रखिए (कैलस);
प्ररोह प्रेरित करने के लिए साइटोकाइनिन-ऑक्सिन अनुपात बढ़ाइए; उन प्ररोहों
को हर कुछ सप्ताह में उपसंवर्धित कीजिए; जड़ लगाने के लिए ऑक्सिन-बहुल
माध्यम पर ले जाइए; और पौधकों को ऊँची आर्द्रता में, फिर मिट्टी में,
अभ्यस्त कीजिए।
\end{solution}

\begin{solution}{exo:b2:asexual-reproduction:5}
हर पौधे की जगह $1 + 24 = 25$ आ जाते हैं: 4 वर्ष बाद $25^{4} =
\num{390625}$ पौधे, जिन्हें \qty{39000}{m^2}, यानी लगभग चार हेक्टेयर चाहिए।
जगह, प्रकाश और पोषक, भूस्तारी की छोटी पहुँच, और भीड़ में पौधों का
मरना उस चरघातांकी प्रावस्था को दो-एक वर्ष में ही रोक देते हैं; और फिर वह
\omterm{def:b2:asexual-reproduction:clone}{क्लोन} एक मोर्चे की तरह आगे बढ़ता है।
\end{solution}

\begin{solution}{exo:b2:asexual-reproduction:6}
तनों की $\num{14000}/130 \approx 110$ पीढ़ियाँ। \qty{4.3e5}{m^2} की
चक्रिका की त्रिज्या: $\sqrt{4.3\times 10^{5}/\pi} = \qty{370}{m}$;
बढ़त प्रति वर्ष $370/\num{14000} = \qty{0.026}{m}$, यानी तीन
सेंटीमीटर से कम।
\end{solution}

\begin{solution}{exo:b2:asexual-reproduction:7}
$p_t = 2^{t}p_0/(1 + (2^{t} - 1)p_0)$: $t = 5$: $0.0032$; $t = 10$:
$0.1024/1.1023 = 0.093$; $t = 15$: $3.28/4.28 = 0.77$। वह कभी नहीं
फैलता यदि उसकी सुयोग्यता लैंगिक वंशक्रम की आधी से कम हो: सापेक्ष
सुयोग्यता $w$ $q' = 2wq$ देती है, जो $w <
\tfrac12$ के लिए क्षीण होती जाती है।
\end{solution}

\begin{solution}{exo:b2:asexual-reproduction:8}
$10 : 1$: जड़ें; $1 : 1$: अविभेदित कैलस; $1 : 10$: प्ररोह।
किसी कलम के पास प्ररोह पहले से है और उसे जड़ें चाहिए, जिन्हें ऑक्सिन
कटे आधार पर प्रेरित करती है।
\end{solution}

\begin{solution}{exo:b2:asexual-reproduction:9}
1, 10, 100, \num{1000}, \num{10000} kg: तीसरी ऋतु के बाद केवल
\qty{1}{t}, और चौथी के बाद \qty{10}{t}: यानी चार ऋतुएँ। असली \omterm{def:b2:angiosperm-reproduction:seed}{बीज}
प्रति पौधा हज़ारों \omterm{def:b2:angiosperm-reproduction:seed}{बीज} देता है, पर आलू \omterm{def:b2:meiosis-heredity:vocabulary}{विषमयुग्मजी}
चतुर्गुणित है, इसलिए अंकुर पृथक हो जाते हैं और कोई भी वह क़िस्म नहीं रहता।
\end{solution}

\begin{solution}{exo:b2:asexual-reproduction:10}
$e^{-5} = 0.0067$: बड़ी आबादी में 670 उत्परिवर्तन-मुक्त व्यक्ति,
और छोटी में लगभग 7। सात व्यक्ति कुछ ही पीढ़ियों में संयोग से सब के सब
संतति छोड़ने में विफल हो सकते हैं (उनमें से किसी के भी जनन न करने की
प्रायिकता प्रति पीढ़ी $e^{-7}$ की कोटि की है), जिसके बाद वह सर्वोत्तम
वर्ग सदा के लिए चला जाता है और दूसरा-सर्वोत्तम सर्वोत्तम बन जाता है:
यानी रैचट का एक खटका, जो बोझ असह्य होने तक दोहराया जाता है। 670 के साथ
वह हानि व्यावहारिक रूप से असंभव है --- जब तक कोई अड़चन $N$ को न घटा दे
या \omterm{thm:b2:genome-diversification:rate}{उत्परिवर्तन दर} न चढ़ जाए; और वह रैचट सदा एक ही दिशा में घूमता है।
\end{solution}

\begin{solution}{exo:b2:asexual-reproduction:11}
उस \omterm{def:b2:asexual-reproduction:clone}{क्लोन} में हर पोषी सुग्राही है: वह विभेद बिना रोक-टोक उसी दर से फैलता
है जिस दर से उसके बीजाणु यात्रा करते हैं। जबकि लैंगिक आबादी में
जिस \omterm{def:b2:meiosis-heredity:vocabulary}{जीन-प्ररूप} को वह हराता है वह $2^{10} = 1024$ संयोजनों में से एक है,
और लगभग हज़ारवें भाग पौधों में उपस्थित है, तथा प्रतिरोधी पड़ोसियों के
बीच बिखरा है, इसलिए वह महामारी भूखी रह जाती है। लंपर एक ही
\omterm{def:b2:asexual-reproduction:clone}{क्लोन} था जो किसी देश के अधिकांश भाग पर उगाया जाता था; उस अंगमारी ने हर
पौधा खुला पाया, और जब तक दूसरी क़िस्में आयात न हुईं तब तक प्रजनन के लिए
कोई विविधता ही नहीं थी।
\end{solution}

\begin{solution}{exo:b2:asexual-reproduction:12}
\omterm{prop:b2:asexual-reproduction:whysex}{मुलर का रैचट}, अनुकूल \omterm{def:b2:genome-diversification:mutation}{उत्परिवर्तनों} का जुड़ना, और
लाल रानी --- हर एक लैंगिकता को ऐसा लाभ देता है जो किसी बड़े,
परजीवी-ग्रस्त, बदलते संसार में दो के गुणक से अधिक हो सकता है। ब्डेलॉइड
रोटिफ़र कठिन उदाहरण हैं: बिना नरों के चार करोड़ वर्ष और कोई रैचट-गलन
नहीं। लगता है वे अत्यधिक निर्जलन-सहिष्णुता से (जो उनके परजीवी मार देती
है), पराया DNA ग्रहण करके और असामान्य रूप से दक्ष \omterm{prop:b2:genome-diversification:repair}{DNA मरम्मत} से बच निकलते
हैं --- यानी ऐसे अपवाद जो दिखाते हैं कि वह नियम सामान्यतः क्या माँगता है।
\end{solution}

\begin{solution}{pb:b2:asexual-reproduction:1}
\textbf{1.} $4\times 10^{4}$ पौधे।
\textbf{2.} $1 + 6\times 2 = 13$।
\textbf{3.} \num{1300}; \num{16900}; \num{219700}।
\textbf{4.} $13^{t} = 400$: $t = \ln 400/\ln 13 = 5.99/2.56 = 2.3$
वर्ष --- यानी तीसरे वर्ष के दौरान भर जाता है।
\textbf{5.} एक हेक्टेयर की चक्रिका की त्रिज्या $\sqrt{10^{4}/\pi} = \qty{56}{m}$;
प्रति वर्ष \qty{0.5}{m} पर \num{113} वर्ष।
\textbf{6.} खेत भर में बिखरे सौ संस्थापक दो वर्षों में सौ मोर्चों से
उसे चरघातांकी रूप से भर देते हैं; जबकि अकेला संस्थापक अपना पड़ोस दो
वर्षों में भरता और फिर सदी भर रेंगता।
\textbf{7.} $5^{n} \ge 10^{6}$: $n \ge 8.6$, यानी नौ उपसंवर्धन
($1.95\times 10^{6}$ प्ररोह), 36~सप्ताह।
\textbf{8.} $0.8\times 1.95\times 10^{6} = 1.6\times 10^{6}$ पौधे।
\textbf{9.} $0.98$; $0.98^{10} = 0.82$।
\textbf{10.} $10\times 0.98 = 9.8$ स्वच्छ वंशक्रम प्रत्याशित; जबकि दस
कलमों में से एक।
\textbf{11.} वह विषाणु जीवद्रव्यतंतुओं और फ्लोएम से यात्रा करता है,
और शीर्षस्थ गुंबद में कोई संवहनी संबंध नहीं है तथा वह विषाणु के फैलने से
तेज़ विभाजित होता है, इसलिए बाहरी दसवाँ मिलिमीटर प्रायः असंक्रमित रहता है।
\textbf{12.} प्रयोगशाला: तुरंत एकरूप, विषाणु-मुक्त पौधे, पर
महँगी, और संवर्धन में कोई भी \omterm{def:b2:genome-diversification:mutation}{उत्परिवर्तन} या संदूषण दस लाख गुना
हो जाता है। भूस्तारी: मुफ़्त और मौक़े पर, पर तीन
ऋतुएँ गईं और मातृ पौधों का हर विषाणु साथ चला।
\textbf{13.} असंगजनक $pN$ हैं और हर एक $F$ \omterm{def:b2:angiosperm-reproduction:seed}{बीज} देता है; लैंगिक
$(1 - p)N$ हैं और हर एक $F/2$ \omterm{def:b2:angiosperm-reproduction:seed}{बीज} देता है (उनका आधा प्रयास \omterm{def:b2:plant-life-cycles:heterospory}{पराग} में
जाता है, जो कोई \omterm{def:b2:angiosperm-reproduction:seed}{बीज} नहीं बनाता); इसलिए \omterm{def:b2:angiosperm-reproduction:seed}{बीजों} में असंगजनक हिस्सा $Fp/(Fp + F(1 - p)/2)
= 2p/(1 + p)$ है।
\textbf{14.} $p_5 = 0.032/1.031 = 0.031$; $p_{10} = 1.024/2.023 =
0.51$: वह पीढ़ी 10 पर एक-आधा पार करता है।
\textbf{15.} $p' = 2p(1 - cp)/\bigl(2p(1 - cp) + (1 - p)\bigr)$।
\textbf{16.} $p' = p$ तब, जब $2(1 - cp^*) = 1$: $p^* = 1/2c = 0.625$।
$p^*$ से नीचे $2(1 - cp) > 1$ है और $p$ चढ़ता है; ऊपर वह गिरता है: यानी स्थिर।
\textbf{17.} $p^* \ge 1$ तब, जब $c \le \tfrac12$: यानी असंगजनक छा जाता
है, जब तक वह परजीवी उसके आम होने पर उसे उसके निर्गत के आधे से अधिक का दाम न लगाए।
\textbf{18.} $c = 0.6$: $p^* = 0.83$। आबादी में लैंगिकता बनाए रखने के
लिए उस परजीवी को, असंगजनक की ऊँची बारंबारता पर, उस \omterm{def:b2:asexual-reproduction:clone}{क्लोन} की
सुयोग्यता दुगुने लाभ से अधिक काटनी पड़ेगी --- यानी लाल रानी को
तेज़ दौड़ना पड़ेगा।
\textbf{19.} तीन गुणसूत्र-समुच्चय \omterm{def:b2:meiosis-heredity:meiosis}{अर्धसूत्री विभाजन} में युग्मित नहीं हो
सकते, इसलिए उसके युग्मक असंतुलित हैं; वह पुनर्संयोजन से बाहर है, और
\omterm{prop:b2:asexual-reproduction:whysex}{मुलर का रैचट} तथा परजीवी उस पर बिना रोक-टोक काम करते हैं --- और इसीलिए
असंगजनक डैंडेलायन वंशक्रम नए हैं और बार-बार लैंगिक पूर्वजों से फिर बनाए
जाते हैं।
\textbf{20.} प्रति विभाजन $10^{-9}\times 4.8\times 10^{8} = 0.48$।
\textbf{21.} $2\times \num{14000}\times 20 = 5.6\times 10^{5}$
विभाजन; $0.48\times 5.6\times 10^{5} = 2.7\times 10^{5}$
\omterm{def:b2:genome-diversification:mutation}{उत्परिवर्तन}।
\textbf{22.} जीनोम का $2.7\times 10^{5}/4.8\times 10^{8} = 5.6\times 10^{-4}$;
यानी लगभग \num{5400} कूटलेखी अंतर।
\textbf{23.} विभज्योतक में कोशिका-वंशक्रमों के बीच होड़ (धीरे विभाजित
होती उत्परिवर्ती पंक्ति हटा दी जाती है) और \omterm{def:b2:asexual-reproduction:clone}{रैमेटों} के बीच (जो तना
कोई बुरा \omterm{def:b2:genome-diversification:mutation}{उत्परिवर्तन} ढोए वह कम मूल-प्ररोह देता है)। कमज़ोर इसलिए कि
अधिकांश कायिक \omterm{def:b2:genome-diversification:mutation}{उत्परिवर्तन} \omterm{def:b2:meiosis-heredity:vocabulary}{विषमयुग्मजी} और ढँके होते हैं, और कुछ भी उन्हें
उजागर नहीं करता: न कोई \omterm{def:b2:meiosis-heredity:meiosis}{अर्धसूत्री विभाजन} उन्हें \omterm{def:b2:meiosis-heredity:vocabulary}{समयुग्मजी} बनाता है, न
कोई एक-कोशिकी युग्मनज अवस्था पूरे जीनोम को एक ही कोशिका में परखती है।
\textbf{24.} एकरूप नहीं: उसके तने लाखों स्थलों पर भिन्न हैं। पर वे
अंतर अधिकतर उदासीन और ढँके हैं, इसलिए वह
\omterm{def:b2:asexual-reproduction:clone}{क्लोन} व्यवहार में एक ही \omterm{def:b2:meiosis-heredity:vocabulary}{जीन-प्ररूप} है --- यानी एक \omterm{def:b2:asexual-reproduction:clone}{जेनेट} के कायिक विभेदों
का मोज़ेक।
\textbf{25.} खेत 2.3 वर्ष बाद भर जाता है; एक विभज्योतक नौ महीनों में
$1.6\times 10^{6}$ पौधे देता है; और $c = 0.8$ के साथ वह असंगजनक
$p^* = 0.625$ पर ठहर जाता है, और लैंगिक \qty{37.5}{\%} रख लेते हैं।
\end{solution}
